\section{Articulación de Actores de Ciencia, Tecnología e Innovación durante la Ejecución de la Propuesta}

El desarrollo de la presente propuesta se llevará a cabo mediante la articulación de actores académicos, clínicos e institucionales, en el marco del Sistema Nacional de Ciencia, Tecnología e Innovación (SNCTeI), con el fin de garantizar un enfoque interdisciplinario que integre capacidades en inteligencia artificial, procesamiento de señales biomédicas y práctica clínica en salud mental, a través de la colaboración activa entre actores con roles complementarios.

Desde el componente académico, el proyecto será desarrollado en el marco del Grupo de Investigación en Automática de la Universidad Tecnológica de Pereira (UTP), clasificado en categoría A1, el nivel más alto otorgado en Colombia, según los resultados de la Convocatoria 957 de 2024 del SNCTeI. Este grupo aportará su experiencia en modelado matemático, aprendizaje automático y procesamiento de señales. Además, el grupo contribuirá en el diseño, implementación y validación de modelos fundacionales y enfoques de inferencia estocástica aplicados al análisis de señales EEG, así como en la formación de recurso humano altamente calificado.

En el ámbito clínico, se contempla la articulación con instituciones especializadas en neurociencias y salud mental de la región, tales como Neurocentro, el Instituto del Sistema Nervioso de Risaralda, el Hospital de Caldas, el Hospital Infantil Universitario Rafael Henao Toro (Manizales) y la Cruz Roja Colombiana, las cuales han mantenido una relación activa con el Grupo de Investigación en Automática a través de la ejecución conjunta de proyectos de investigación financiados por MINCIENCIAS. Estas entidades cuentan con experiencia en diagnóstico, tratamiento y seguimiento de trastornos neurológicos y psiquiátricos, así como con infraestructura para la adquisición de señales EEG y evaluación clínica especializada. Neurocentro, por ejemplo, se posiciona como un referente regional en neurociencias con un enfoque integral e interdisciplinario, mientras que el Instituto del Sistema Nervioso de Risaralda dispone de servicios especializados en salud mental y hospitalización, lo que facilita la generación de escenarios reales para la validación clínica de las metodologías propuestas.

Adicionalmente, se promoverá la vinculación con organizaciones y asociaciones regionales orientadas al estudio y atención de trastornos neurológicos y del neurodesarrollo, tales como la Fundación Neuropsikids, enfocada en el abordaje integral del neurodesarrollo infantil, la Asociación Colombiana de Neuromodulación, orientada a la investigación y aplicación de técnicas terapéuticas en neurociencias, y la Asociación Sinapsis, entre otras iniciativas locales relacionadas con salud mental y procesos neurológicos. Estas organizaciones permitirán fortalecer los procesos de apropiación social del conocimiento, facilitar la divulgación científica y promover la interacción con la comunidad.

Esta articulación permitirá no solo el acceso a datos clínicos y escenarios reales de aplicación, sino también la retroalimentación continua por parte de especialistas en salud, garantizando la pertinencia clínica de las soluciones desarrolladas. Asimismo, favorecerá la validación interdisciplinaria de los resultados y la transferencia del conocimiento hacia el entorno clínico y social, contribuyendo al fortalecimiento del ecosistema regional de ciencia, tecnología e innovación en el área de la salud mental.